\selectlanguage{ngerman}%
\frenchspacing%
\begin{ledgroupsized}[r]{120mm}%
\footnotesize%
\pstart%
\noindent\textbf{Überlieferung:}%
\pend%
\end{ledgroupsized}%
\begin{ledgroupsized}[r]{114mm}%
\footnotesize%
\pstart%
\parindent -6mm%
\makebox[6mm][l]{\textit{L}}%
Aufzeichnung:
LH~XXXVII,\;5 Bl.\;171--172.
Ein Bogen, 8\textsuperscript{o},
Ränder zum Teil beschnitten;
Papiererhaltungsmaßnahmen.
\href{\Scan{LH_37_05_171-172wz}}{Fragment} eines Wasserzeichens am oberen Rand.
Eine ganzseitig beschriebene Seite auf Bl.\;\href{\Scan{LH_37_05_172v}}{172\,v\textsuperscript{o}}.
Bl.\;\href{\Scan{LH_37_05_171r}}{171\,r\textsuperscript{o}} überliefert die Aufzeichnung N.\;\ref{60349-1}.
Bl.\;\href{\Scan{LH_37_05_171v}}{171\,v\textsuperscript{o}} und
Bl.\;\href{\Scan{LH_37_05_172r}}{172\,r\textsuperscript{o}} sind unbeschrieben.
\pend%
\end{ledgroupsized}%
\vspace{5mm}%
\begin{ledgroup}%
\footnotesize%
\pstart%
\noindent%
\textbf{Datierungsgründe:}
Die titellose und undatierte Aufzeichnung N.\;\ref{60349-2}, in der Leibniz vor dem Hintegrund einer kritischen Stellungnahme zu Descartes' Lehre der Elemente nach der Natur des Lichts fragt, ist aller Wahrscheinlichkeit nach zur gleichen Zeit entstanden wie die auf demselben Trager überlieferte Aufzeichnung N.\;\ref{60349-1}.
Siehe zur Datierung dieser letzteren die ausführliche Begründung auf S.\;\pageref{60349-1_dat}.
\pend%
\end{ledgroup}%
\selectlanguage{latin}%
\frenchspacing%
\vspace{8mm}%
\normalsize%
\pstart%
\noindent%
\bracks{%
\href{\Scan{LH_37_05_172v}}%
{172\,v\textsuperscript{o}}%
}
\edtext{%
Cartesius\indexN{\textls{Descartes} (Cartesius, des Cartes), René 1596\textendash1650}
in suo libro%
\indexS{liber}
quem \textit{mundum} inscripsit,
sibi persuadet
elementa sua pura%
\indexS{elementa Cartesii}%
\indexS{elementum purum}
esse in maxima universi parte,%
\indexS{universum}
mixta%
\indexS{corpus mixtum}%
\indexS{mixta}
autem tantum in locis %
\indexS{locus}
ubi sibi vicina sunt%
\bracks{;}
ita materiam subtilem%
\indexS{materia subtilis}
seu fluidum perfectum%
\indexS{fluidum perfectum}
collocat in fixis,%
\indexS{stella fixa}
\edtext{%
globulos secundi elementi%
\indexS{globulus secundi elementi}%
\indexS{elementum secundum}%
}{%
\lemma{globulos}%
\Bfootnote{%
\textit{(1)}~aut \angles{--} %
\textit{(2)}~secundi elementi%
~\textit{L}%
}}
credit implere coelos;%
\indexS{coelum}
tertium elementum,%
\indexS{elementum tertium}
quod nec divisibilitatem primi,%
\indexS{divisibilitas primi elementi}%
\indexS{elementum primum}
nec mobilitatem secundi elementi%
\indexS{mobilitas secundi elementi}
habet%
\bracks{,}
statuit locum habere
in planetis%
\indexS{planeta}
et cometis.%
\indexS{cometa}
A primo elemento%
\indexS{elementum primum}
esse lucis originem,%
\indexS{origo lucis}%
\indexS{lux}
a secundo%
\indexS{elementum secundum}
ejus transmissionem,%
\indexS{transmissio lucis}
a
\edtext{%
tertio%
\indexS{elementum tertium}
resistentiam%
\indexS{resistentia et opacitas}
et opacitatem%
\indexS{opacitas}.%
}{%
\lemma{tertio}%
\Bfootnote{%
\textit{(1)}~opacitatem seu %
\textit{(2)}~\angles{reflex} %
\textit{(3)}~resistentiam et opacitatem.%
~\textit{L}%
}}%
}{%
\lemma{Cartesius \bracks{...} opacitatem}%
\Cfootnote{%
Vgl.
\textsc{R.\;Descartes,} \textit{Le monde}, cap.\;V
(Paris 1664, S.\;59--64;\cite{02532}
\textit{DO} XI, S.\;28.14--30.27).\cite{00120}%
}}
\pend%
\pstart%
Sed mihi parum divinae magnitudini%
\indexS{magnitudo divina}
consentaneum videtur,
\edtext{%
ut varietas%
\indexS{varietas creaturarum}%
}{%
\lemma{ut}%
\Bfootnote{%
\textit{(1)}~varietati %
\textit{(2)}~varietas%
~\textit{L}%
}}
creaturarum%
\indexS{creatura}
in solis locum habeat
\edtext{%
confiniis illis%
\indexS{confinium}%
}{%
\lemma{confiniis}%
\Bfootnote{%
\textit{(1)}~\angles{rerum,} %
\textit{(2)}~illis,%
~\textit{L}%
}}
quae respectu ingentium illarum massarum%
\indexS{massa ingens}
nullius momenti sunt,
et pene punctorum instar.
Quin contra arbitror
nullam esse particulam universi,%
\indexS{universum}%
\indexS{particula universi}
in qua non existat
quidam creaturarum%
\indexS{creatura}
velut Mundus.%
\indexS{mundus}
\pend%
\pstart%
Deinde
\edtext{%
quod Lucem considerat%
\indexS{lux}
velut elementum quoddam%
\indexS{elementum resolutionis ultimae}%
\indexS{elementum subtilitatis extremae}
ultimae resolutionis%
\indexS{resolutio ultima}
extremaeque subtilitatis,%
\indexS{subtilitas extrema}%
}{%
\lemma{quod}%
\Bfootnote{%
\textit{(1)}~Lumen %
\textit{(2)}~Lucem considerat %
\textit{(a)}~quasi primum quoddam %
\textit{(b)}~velut elementum quoddam %
\textit{(aa)}~extremae %
\textit{(bb)}~ultimae resolutionis extremaeque subtilitatis,%
~\textit{L}%
}%
\lemma{Lucem \bracks{...} subtilitatis}%
\Cfootnote{%
Siehe a.a.o. (S.\;62;\cite{02532}
\textit{DO} XI, S.\;29.28--30.2).\cite{00120}
}%
}
in eo assentiri non
\edtext{%
possum.
Nam cum phosphoros%
\indexS{phosphorus Bononiensis}
considero,
maxime
\edtext{%
Bononiensem%
\indexS{phosphorus Bononiensis}
qui lucem ex aere attrahit,%
\indexS{lux attractus ex aere}%
\indexS{aer}%
}{%
\lemma{Bononiensem \bracks{...} attrahit}%
\Cfootnote{%
Gemeint ist der phos­pho­res­zie­rende \textit{lapis Bononiensis}.
Den Bologneser Leuchtstein kannte Leibniz etwa aus
\textsc{H.\;Fabri}, \textit{Physica},
tract.\;II, lib.\;III, prop.\;64
(Lyon 1669--1671, Bd.\;I, S.\;442b--451a)\cite{00044}
und
\textsc{O.\;v.\;Guericke}, \textit{Experimenta nova},
lib.\;IV, cap.\;10 (Amsterdam 1672, S.\;139b);\cite{00055}
beide Stellen hatte er 1670--1672 exzerpiert
(vgl. \textit{LSB} VIII,\;2 N.\;55, S.\;499.20--501.21\cite{01234}
bzw. \textit{LSB} VIII,\;1 N.\;36, S.\;251.18\,f.).\cite{01197}
Eine weitere Quelle war möglicherweise:
\textsc{A.\;Kircher}, \textit{Magnes},
l.\;III, p.\;II, c.\;4, q.\;2 (Rom 1654, S.\;463).\cite{02533}%
}}
\ledpb%
suspicor
lumen%
\indexS{lumen}%
}{%
\lemma{possum.}%
\Bfootnote{%
\textit{(1)}~Credo %
\textit{(2)}~Nam cum phosphoros considero, %
\textit{(a)}~lucemque %
\textit{(b)}~maxime Bononiensem %
\bracks{...} aere attrahit, %
\textit{(aa)}~in f\angles{ac--} %
\textit{(bb)}~suspicor lumen%
~\textit{L}%
}}
esse liquorem aliquem%
\indexS{liquor subtilitatis ultimae}
non utique ultimae subtilitatis,%
\indexS{subtilitas ultima}
alioqui non ita facile a corpore hujusmodi hauriretur;%
\indexS{liquor haustus a corpore}%
\indexS{corpus hauriens liquorem}
nam ipsum ultimae subtilitatis%
\indexS{subtilitas ultima}
fluidum%
\indexS{fluidum subtilitatis ultimae}
corpore aliquo crasso%
\indexS{corpus crassum}
et haud dubie poroso%
\indexS{corpus porosum}
capere velle,
foret aquam haurire cribro.%
\indexS{haustus aquae cribro}%
\indexS{cribrum}
Itaque
\edtext{%
nec instantaneam esse arbitror luminis transmissionem,%
\indexS{transmissio luminis instantanea}
etsi non negem %
liquidum%
}{%
\lemma{nec}%
\Bfootnote{%
\textit{(1)}~instantaneum %
\textit{(2)}~instantaneam esse arbitror luminis %
\textit{(a)}~motum %
\textit{(b)}~transmissionem, etsi %
\textit{(aa)}~liquidum %
\textit{(bb)}~non negem, liquidum%
~\textit{L}%
}}
quo transmittitur%
\indexS{liquidum quo lumen transmittitur}
ab astris%
\indexS{aster}
ad nos usque
diffusum esse.%
\indexS{liquidum diffusum}
Nec simplici impulsu,%
\indexS{impulsus}
sed certa quadam modificatione motus,%
\indexS{modificatio motus}
qualem et in sono%
\indexS{sonus}
agnoscimus,
credo propagari%
\indexS{propagatio luminis}
atque ita liquorem illum%
\indexS{liquor accensus per aerem}
ut ita dicam
per totum aerem%
\indexS{aer}
accendi%
\indexS{liquor accensus per aerem}
instar fumi%
\indexS{fumus}
aut
\edtext{%
fulguris,%
\indexS{fulgur}
aut halitus naphthae,%
\indexS{halitus naphthae}%
\indexS{naphtha}
licet incomparabili subtilitate%
\indexS{subtilitas incomparabilis}%
}{%
\lemma{fulguris, aut}%
\Bfootnote{%
\textit{(1)}~subtilissimi %
\textit{(2)}~halitus naphthae, licet incomparabili subtilitate,%
~\textit{L}%
}}
et velocitate.%
\indexS{velocitas incomparabilis}
Et peculiarem accensionis%
\indexS{accensio liquidi per aerem}
sive motus explodentis modum%
\indexS{modus motus}%
\indexS{motus explodens}
esse colorem.%
\indexS{color}
\pend%
